\documentclass[11pt,a4paper]{article}

% ------------------------------------------------------------------
% Revised proofs of the EIT No-Go/Go theory
% Version 7.0, 2026-08-06
%
% This file replaces the asymptotic core of the v6.2 proof package.
% The main asymptotic theorem is proved by the Laurent expansion of a
% finite-dimensional rational resolvent.  The Neumann expansion is retained
% as a quantitative lemma for a computable convergence domain and remainder.
% The former sector-graph remark is promoted to a theorem.
% ------------------------------------------------------------------

\usepackage[margin=25mm]{geometry}
\usepackage{amsmath,amssymb,amsthm,mathtools}
\usepackage{bm}
\usepackage{booktabs}
\usepackage{enumitem}
\usepackage[colorlinks=true,linkcolor=blue!60!black]{hyperref}
\usepackage{tcolorbox}
\tcbuselibrary{skins,breakable}

\newtheorem{theorem}{Theorem}[section]
\newtheorem{lemma}[theorem]{Lemma}
\newtheorem{corollary}[theorem]{Corollary}
\newtheorem{proposition}[theorem]{Proposition}
\newtheorem{definition}[theorem]{Definition}
\theoremstyle{remark}
\newtheorem{remark}[theorem]{Remark}

\newtcolorbox{scopebox}[1][Scope]{
  breakable,enhanced,colback=yellow!5,colframe=yellow!45!black,
  title=#1,fonttitle=\bfseries}
\newtcolorbox{claimbox}[1][Claim discipline]{
  breakable,enhanced,colback=blue!4,colframe=blue!50!black,
  title=#1,fonttitle=\bfseries}

\newcommand{\im}{\mathrm{i}}
\newcommand{\Tr}{\operatorname{Tr}}
\newcommand{\rank}{\operatorname{rank}}
\newcommand{\adj}{\operatorname{adj}}
\newcommand{\spann}{\operatorname{span}}
\newcommand{\spec}{\operatorname{spec}}
\newcommand{\Ran}{\operatorname{Ran}}
\newcommand{\dist}{\operatorname{dist}}
\newcommand{\ket}[1]{|#1\rangle}
\newcommand{\bra}[1]{\langle #1|}
\newcommand{\braket}[2]{\langle #1|#2\rangle}
\newcommand{\dchi}{\delta\chi_S}

\title{\bfseries Revised Proofs of the EIT No-Go/Go Theory\\[4pt]
\large Laurent asymptotics, sector-graph selection, and observable inheritance}
\author{Quantum Sensing Group}
\date{August 6, 2026}

\begin{document}
\maketitle

\begin{abstract}
We give a finite-dimensional proof package for the no-go/go theory used in the
NV--EIT and companion PRL programmes.  The large-dissipation expansion is
proved from the Laurent expansion at infinity of a rational resolvent; no
operator-norm hypothesis is needed for the asymptotic order theorem.  A
Neumann-series lemma is retained because it supplies a directly checkable
convergence condition and an explicit truncation error.  We promote the
sector-graph selection rule from a remark to a theorem, carefully separating a
graph-theoretic lower bound from equality, which additionally requires a
nonvanishing sum of shortest-path amplitudes.  We also close the omitted
algebra in the pure stationary-state theorem, remove the Neumann dependence
from the exact-transfer-zero theorem, state the observable-order inheritance
rule, and make the Hermitian and semisimple restrictions of the singular
results explicit.
\end{abstract}

\tableofcontents

% ==================================================================
\section{Scope, notation, and the three meanings of no-go}
% ==================================================================

\begin{scopebox}
The results below assume a finite-dimensional, Markovian, weak-probe linear
response problem.  The fixed-scaling family is
\[
A_\Gamma(z)=\Gamma D+A_0(z),
\]
where $\Gamma\to\infty$ is the single large dissipation scale and $z$ denotes
frequency or other parameters held fixed.  A second scale that itself grows
with $\Gamma$ is a different two-scale problem and is not silently included in
these theorems.
\end{scopebox}

For an invertible $D$, define
\begin{equation}
X(z):=-D^{-1}A_0(z),\qquad \nu:=D^{-1}c,
\qquad K_\Gamma(z):=p^\dagger A_\Gamma(z)^{-1}c.
\label{eq:standing}
\end{equation}
Then
\begin{equation}
A_\Gamma=D(\Gamma I-X),
\qquad
K_\Gamma=p^\dagger(\Gamma I-X)^{-1}\nu.
\label{eq:factor}
\end{equation}
The path moments are
\begin{equation}
M_n(z):=p^\dagger X(z)^n\nu,
\qquad n=0,1,2,\ldots.
\label{eq:moments}
\end{equation}

We distinguish three statements that must not be conflated.
\begin{enumerate}[label=(\roman*)]
\item An \emph{exact no-go} is $K_\Gamma\equiv0$ on its regular domain.
\item An \emph{asymptotic no-go/order statement} is
      $K_\Gamma=O(\Gamma^{-r})$ as $\Gamma\to\infty$.
\item An \emph{experimental no-go} says that the engineerable, asymptotic,
      and detectable parameter windows have empty intersection.
\end{enumerate}
Only (i) is a transfer zero.  Statement (ii) permits a nonzero finite-$\Gamma$
response, and statement (iii) is a property of a measurement budget, not an
algebraic identity.

% ==================================================================
\section{Optical dark vectors and stationary pure dark states}
% ==================================================================

\begin{theorem}[Optical dark-subspace rank]
Let $\Omega:\mathcal H_g\to\mathcal H_e$ be the optical coupling map.  The
optical dark subspace is $\ker\Omega$, and
\[
\dim\ker\Omega=\dim\mathcal H_g-\rank\Omega.
\]
The dimension is invariant under invertible basis changes on either manifold.
\end{theorem}

\begin{proof}
This is rank--nullity.  Left or right multiplication by an invertible matrix
composes $\Omega$ with a bijection and therefore does not change its rank.
\end{proof}

Let
\begin{equation}
\mathcal L(\rho)=-\im[H,\rho]
+\sum_\mu\left(L_\mu\rho L_\mu^\dagger
-\frac12\{L_\mu^\dagger L_\mu,\rho\}\right).
\label{eq:gksl}
\end{equation}

\begin{lemma}[Purity-variance identity]
For $\rho_D=\ket D\bra D$ with $\braket DD=1$,
\[
\frac{d}{dt}\Tr\rho^2\Big|_{\rho_D}
=-2\sum_\mu\left(
\bra D L_\mu^\dagger L_\mu\ket D
-|\bra D L_\mu\ket D|^2\right)\le0.
\]
Equality holds iff every $L_\mu\ket D$ is parallel to $\ket D$.
\end{lemma}

\begin{proof}
Use $\frac{d}{dt}\Tr\rho^2=2\Tr(\rho\mathcal L\rho)$, cyclicity of the trace,
and $\rho_D^2=\rho_D$.  The Hamiltonian contribution vanishes, while each
jump contributes minus twice the squared norm of
$(L_\mu-\bra D L_\mu\ket D)\ket D$.
\end{proof}

\begin{theorem}[Stationary pure Lindblad state: corrected proof]
The rank-one state $\rho_D=\ket D\bra D$ is stationary iff there exist
$\lambda_\mu\in\mathbb C$ and $h\in\mathbb R$ such that
\begin{align}
L_\mu\ket D&=\lambda_\mu\ket D\quad(\forall\mu),
\label{eq:jump-eig}\\
\widetilde H\ket D&=h\ket D,
\qquad
\widetilde H:=H+\frac{\im}{2}\sum_\mu
(\lambda_\mu^*L_\mu-\lambda_\mu L_\mu^\dagger).
\label{eq:Htilde}
\end{align}
\end{theorem}

\begin{proof}
If $\rho_D$ is stationary, its purity is constant.  The preceding lemma gives
\eqref{eq:jump-eig}.  For one jump satisfying $L\ket D=\lambda\ket D$, set
\[
K_\lambda:=\frac{\im}{2}(\lambda^*L-\lambda L^\dagger)=K_\lambda^\dagger.
\]
Using only $L\rho_D=\lambda\rho_D$ and
$\rho_D L^\dagger=\lambda^*\rho_D$, direct multiplication gives
\begin{align*}
-\im[K_\lambda,\rho_D]
&=|\lambda|^2\rho_D
-\frac{\lambda}{2}L^\dagger\rho_D
-\frac{\lambda^*}{2}\rho_D L\\
&=L\rho_D L^\dagger
-\frac12\{L^\dagger L,\rho_D\}.
\end{align*}
Thus, after summing over jumps,
\[
\mathcal L(\rho_D)=-\im[\widetilde H,\rho_D].
\]
Stationarity is therefore equivalent to
$[\widetilde H,\rho_D]=0$.  A Hermitian operator commutes with a rank-one
projector iff its range vector is an eigenvector, yielding
\eqref{eq:Htilde} with real $h$.  Reversing the calculation proves
sufficiency.  Notice that the proof never assumes the generally false
identity $\rho_D L=\lambda^*\rho_D$.
\end{proof}

\begin{remark}
An optical dark vector need not be a stationary pure state.  Ground-state
dephasing can preserve $\Omega\ket D=0$ while violating
$L_\mu\ket D\parallel\ket D$.
\end{remark}

% ==================================================================
\section{The operational no-go object: Schur-complement response}
% ==================================================================

Let the weak-probe equations be
\[
\begin{pmatrix}A&B\\C&G_g\end{pmatrix}
\begin{pmatrix}x\\s\end{pmatrix}
=\begin{pmatrix}b_p\\0\end{pmatrix},
\]
and let the detector act only on $x$.  The pathway-cut counterfactual removes
$B$ and $C$ while keeping all other parameters fixed.

\begin{theorem}[Sector-resolved response formula]
Assume $A$ and $S_g:=G_g-CA^{-1}B$ are invertible.  Then
\begin{equation}
\dchi:=\chi_{\rm full}-\chi_{\rm cut}^{(S)}
=\chi_0 d^\dagger A^{-1}B\,S_g^{-1}C A^{-1}b_p.
\label{eq:schur-response}
\end{equation}
\end{theorem}

\begin{proof}
From $Ax+Bs=b_p$,
$x=A^{-1}b_p-A^{-1}Bs$.  Substitution into the second row gives
$S_gs=-CA^{-1}b_p$, and hence
$s=-S_g^{-1}CA^{-1}b_p$.  Substitute this back into $x$ and subtract the
cut response $\chi_0d^\dagger A^{-1}b_p$.
\end{proof}

\begin{claimbox}
The exact no-go object is the full-minus-cut response $\dchi$, not the total
response $\chi_{\rm full}$.  Perfect EIT can have $\chi_{\rm full}=0$ while
$\dchi\ne0$.  Therefore a zero of the total susceptibility is not by itself a
no-go certificate.
\end{claimbox}

% ==================================================================
\section{Laurent expansion at infinity and the first nonzero moment}
% ==================================================================

\begin{lemma}[Neumann expansion with explicit error]
\label{lem:neumann-v7}
For any submultiplicative norm, if
$q:=\|X\|/|\Gamma|<1$, then
\begin{equation}
K_\Gamma=\sum_{n=0}^{m}\frac{M_n}{\Gamma^{n+1}}+R_m(\Gamma),
\label{eq:neumann-trunc}
\end{equation}
with
\begin{equation}
|R_m(\Gamma)|\le
\frac{\|p\|\,\|D^{-1}\|\,\|c\|}{|\Gamma|}
\frac{q^{m+1}}{1-q}.
\label{eq:neumann-error}
\end{equation}
\end{lemma}

\begin{proof}
By \eqref{eq:factor},
\[
(\Gamma I-X)^{-1}=\Gamma^{-1}(I-X/\Gamma)^{-1}.
\]
Apply the matrix geometric series and its standard tail bound, then contract
with $p^\dagger$ and $\nu=D^{-1}c$.
\end{proof}

\begin{remark}
The Neumann condition is a convenient sufficient condition, not the logical
basis of the asymptotic order theorem.  For a nonnormal $X$, $\|X\|$ can be
much larger than its spectral radius, so the norm condition may be very
conservative.
\end{remark}

\begin{theorem}[Laurent moment expansion]
\label{thm:laurent}
For fixed finite-dimensional $X$, the scalar transfer function
\[
K_\Gamma=p^\dagger(\Gamma I-X)^{-1}\nu
\]
is rational in $\Gamma$ and has a convergent Laurent expansion at infinity,
\begin{equation}
K_\Gamma=\sum_{n=0}^{\infty}\frac{M_n}{\Gamma^{n+1}},
\qquad M_n=p^\dagger X^n\nu,
\label{eq:laurent}
\end{equation}
for all sufficiently large $|\Gamma|$.  Its radius in the variable
$w=1/\Gamma$ is the distance from $w=0$ to the nearest pole of
$(I-wX)^{-1}$; no operator-norm hypothesis is required for the existence of
the expansion.
\end{theorem}

\begin{proof}
Cramer's rule gives
\[
K_\Gamma=
\frac{p^\dagger\adj(\Gamma I-X)\nu}{\det(\Gamma I-X)},
\]
so $K_\Gamma$ is rational.  Put $w=1/\Gamma$.  Then
\[
K_{1/w}=w\,p^\dagger(I-wX)^{-1}\nu.
\]
Because $\det(I-wX)=1+O(w)$, the matrix $(I-wX)^{-1}$ is analytic at
$w=0$.  Let its Taylor series be $\sum_{n\ge0}C_nw^n$.  Multiplying by
$I-wX$ and comparing powers gives $C_0=I$ and $C_n=XC_{n-1}=X^n$.
Therefore
\[
K_{1/w}=\sum_{n\ge0}(p^\dagger X^n\nu)w^{n+1},
\]
which is \eqref{eq:laurent} after returning to $\Gamma$.
\end{proof}

\begin{corollary}[First nonzero moment selects the order]
\label{cor:first-moment}
If
\[
M_0=\cdots=M_{m-1}=0,
\qquad M_m\ne0,
\]
then
\begin{equation}
K_\Gamma=M_m\Gamma^{-m-1}+O(\Gamma^{-m-2}).
\label{eq:first-order}
\end{equation}
Thus the amplitude order is $m+1$.
\end{corollary}

\begin{proof}
Delete the vanishing coefficients from the Laurent series in
Theorem~\ref{thm:laurent}.  The next coefficient is nonzero.
\end{proof}

\begin{proposition}[Uniform remainder on compact parameter sets]
Let $X(z)$, $p(z)$, and $c(z)$ be continuous on a compact set $C$, with
$D(z)$ invertible there.  Then there is a finite
$Q_C:=\sup_{z\in C}\|X(z)\|$.  For $|\Gamma|>Q_C$, the Neumann expansion and
remainder bound are uniform in $z\in C$ after replacing the individual norms
in \eqref{eq:neumann-error} by their suprema on $C$.
\end{proposition}

\begin{proof}
Continuous functions on a compact set are bounded.  Apply
Lemma~\ref{lem:neumann-v7} with the common ratio $Q_C/|\Gamma|<1$.
\end{proof}

% ==================================================================
\section{The sector-graph selection theorem}
% ==================================================================

Let $\{P_\alpha\}_{\alpha\in\mathcal V}$ be mutually orthogonal projectors
with $\sum_\alpha P_\alpha=I$.  Suppose
\begin{equation}
D=\sum_\alpha P_\alpha D P_\alpha,
\qquad
A_0=A_{\rm diag}+V,
\qquad
A_{\rm diag}=\sum_\alpha P_\alpha A_{\rm diag}P_\alpha.
\label{eq:sector-decomp}
\end{equation}
Because $D$ is invertible and block diagonal, so is $D^{-1}$.  Define
\begin{equation}
X_0:=-D^{-1}A_{\rm diag},
\qquad
W:=-D^{-1}V,
\qquad X=X_0+W.
\label{eq:X0W}
\end{equation}
Construct a directed graph $G$ on the sector labels by drawing
$\alpha\to\beta$ iff
\begin{equation}
P_\beta W P_\alpha\ne0.
\label{eq:edge}
\end{equation}

\begin{theorem}[Sector-graph lower bound]
\label{thm:sector-graph}
Assume the source and readout are sector supported,
\[
P_s\nu=\nu,
\qquad
p^\dagger P_t=p^\dagger.
\]
Let $d=\dist_G(s,t)$ be the length of the shortest directed path from $s$ to
$t$, with $d=\infty$ if no path exists.  Then
\begin{equation}
M_n=p^\dagger X^n\nu=0
\qquad\text{for every }n<d.
\label{eq:graph-moment-zero}
\end{equation}
Consequently, if $d<\infty$,
\begin{equation}
K_\Gamma=O(\Gamma^{-d-1}).
\label{eq:graph-order-bound}
\end{equation}
If $d=\infty$, then every moment vanishes and the transfer is exactly zero at
all regular $\Gamma$.
\end{theorem}

\begin{proof}
Expand
\[
X^n=(X_0+W)^n
\]
as a sum of noncommutative words of length $n$ in $X_0$ and $W$.  The block
diagonal operator $X_0$ never changes sector.  Each factor $W$ can change the
sector only along one directed edge of $G$.  A word containing $r$ factors of
$W$ can therefore connect $s$ to $t$ only if there is a directed path of
length at most $r$.  Every word of total length $n<d$ has $r\le n<d$, so its
$t$--$s$ block vanishes.  Contracting with $p^\dagger P_t$ and $P_s\nu$
gives \eqref{eq:graph-moment-zero}.  Equation
\eqref{eq:graph-order-bound} follows from
Corollary~\ref{cor:first-moment}, allowing for additional cancellations.  If
$d=\infty$, the same word argument applies for every $n$; exact vanishing then
follows from Theorem~\ref{thm:krylov} below.
\end{proof}

\begin{corollary}[When graph distance equals the dissipation order]
\label{cor:graph-equality}
Under the assumptions of Theorem~\ref{thm:sector-graph}, suppose $d<\infty$
and the sum of all length-$d$ word amplitudes is nonzero, equivalently
\begin{equation}
M_d=p^\dagger(X_0+W)^d\nu\ne0.
\label{eq:Md-nonzero}
\end{equation}
Then
\begin{equation}
K_\Gamma=M_d\Gamma^{-d-1}+O(\Gamma^{-d-2}).
\end{equation}
\end{corollary}

\begin{remark}[Graph distance is a lower bound, not an automatic equality]
Different shortest paths can interfere and cancel in $M_d$.  In that case the
first nonzero moment occurs at an order larger than $d$.  Therefore the graph
determines the minimum possible number of sector-changing insertions, while
the complex path amplitudes determine whether that minimum survives.  This is
the precise form of the dissipative selection rule used in the PRL claim.
\end{remark}

\begin{remark}[Multiple source or detector sectors]
For sources or detectors supported on several sectors, apply the theorem to
each nonzero pair $(s,t)$.  The earliest potentially nonzero moment is the
minimum graph distance over supported pairs, followed by the same amplitude
cancellation test.
\end{remark}

% ==================================================================
\section{Exact transfer zero without Neumann-series dependence}
% ==================================================================

\begin{definition}
The reachable Krylov space is
\[
\mathcal K(X,\nu):=\spann\{\nu,X\nu,X^2\nu,\ldots\}.
\]
Let $r:=\dim\mathcal K(X,\nu)$.  Equivalently, $r$ is the degree of the
minimal polynomial of $X$ relative to the cyclic vector $\nu$.
\end{definition}

\begin{theorem}[Exact transfer-zero/Krylov criterion]
\label{thm:krylov}
The following are equivalent:
\begin{enumerate}[label=(\roman*)]
\item $M_n=p^\dagger X^n\nu=0$ for $n=0,\ldots,r-1$;
\item $p\perp\mathcal K(X,\nu)$;
\item $p^\dagger q(X)\nu=0$ for every polynomial $q$;
\item $p^\dagger(\Gamma I-X)^{-1}\nu=0$ for every
      $\Gamma\notin\spec X$.
\end{enumerate}
\end{theorem}

\begin{proof}
Conditions (i) and (ii) are the same statement in a Krylov basis.  By
definition of $r$, the vectors $X^n\nu$ for $n\ge r$ are linear combinations
of the first $r$ Krylov vectors, so (i) implies that every moment vanishes and
hence (iii).

For regular $\Gamma$, the adjugate identity writes
$(\Gamma I-X)^{-1}$ as a scalar denominator times a matrix polynomial in $X$;
therefore (iii) implies (iv).

Conversely, the scalar in (iv) is a rational function that vanishes on its
regular domain.  It is therefore the zero rational function.  Its Laurent
expansion at infinity, furnished by Theorem~\ref{thm:laurent}, has every
coefficient zero.  Hence all $M_n$ vanish, in particular the first $r$.
This direction does not invoke a Neumann norm condition.
\end{proof}

\begin{remark}
The older dimension-$N$ test remains correct but is not sharp.  Only $r$
Krylov moments are needed.  Numerically, $r$ is the rank at which the Krylov
sequence stabilizes.
\end{remark}

% ==================================================================
\section{Symmetry-protected exact zeros}
% ==================================================================

\begin{theorem}[Commuting-sector transfer zero]
Let $\{P_\alpha\}$ be orthogonal projectors summing to $I$, and suppose
$[A(z),P_\alpha]=0$ for every $\alpha$.  If
$P_\beta c=c$ and $p^\dagger P_\alpha=p^\dagger$ with $\alpha\ne\beta$,
then
\[
p^\dagger A(z)^{-1}c=0
\]
at every regular point.
\end{theorem}

\begin{proof}
Commutation passes to the inverse.  Therefore
\[
p^\dagger A^{-1}c
=p^\dagger P_\alpha A^{-1}P_\beta c
=p^\dagger A^{-1}P_\alpha P_\beta c=0.
\]
\end{proof}

\begin{remark}[Non-Abelian formulation]
For a non-Abelian symmetry, replace one-dimensional sector labels by the
isotypic projectors of the group representation.  Schur's lemma then gives the
same conclusion when source and detector transform in inequivalent irreducible
sectors.  This is the appropriate language for point groups such as
$C_{3v}$; the commuting-projector theorem is its elementary finite-dimensional
form.
\end{remark}

% ==================================================================
\section{Observable-order inheritance}
% ==================================================================

\begin{lemma}[Arithmetic of nonzero asymptotic factors]
Suppose
\[
f(\Gamma)=a\Gamma^{-r}(1+O(\Gamma^{-1})),\quad
g(\Gamma)=b\Gamma^{-q}(1+O(\Gamma^{-1})),
\]
and
\[
h(\Gamma)=c\Gamma^{-s}(1+O(\Gamma^{-1})),
\qquad abc\ne0.
\]
Then
\[
f(\Gamma)h(\Gamma)^{-1}g(\Gamma)
=\frac{ab}{c}\Gamma^{-(r+q-s)}(1+O(\Gamma^{-1})).
\]
\end{lemma}

\begin{proof}
Invert the nonzero leading expansion of $h$ and multiply the three series.
\end{proof}

\begin{theorem}[Sector-observable inheritance]
\label{thm:inheritance}
For the scalar Schur-complement response
\[
\dchi=\chi_0 L(\Gamma)S_g(\Gamma)^{-1}R(\Gamma),
\]
assume
\[
L\sim a\Gamma^{-n_L},\qquad
R\sim b\Gamma^{-n_R},\qquad
S_g\sim c\Gamma^{-n_S},
\qquad abc\ne0.
\]
Then
\begin{equation}
\dchi\sim\chi_0\frac{ab}{c}
\Gamma^{-\nu_{\rm obs}},
\qquad
\nu_{\rm obs}=n_L+n_R-n_S.
\label{eq:inheritance}
\end{equation}
If the displayed leading coefficient vanishes because of a further matrix or
path cancellation, the actual observable order is higher and must be found
from the first nonzero coefficient of the full product.
\end{theorem}

\begin{remark}[Normalization changes the exponent]
If a reported contrast is $C=\dchi/\chi_{\rm ref}$ and
$\chi_{\rm ref}\sim d\Gamma^{-n_{\rm ref}}$, then
$C$ has order $\nu_{\rm obs}-n_{\rm ref}$, provided $d\ne0$.  Thus signed
absorption difference and normalized contrast need not have the same integer
order.  The observable must be named before quoting an exponent.
\end{remark}

% ==================================================================
\section{Singular damping: proved scope and nonnormal caveat}
% ==================================================================

\begin{theorem}[Hermitian singular damping]
\label{thm:singular-D}
Let $D=D^\dagger\ge0$, let $P$ be the orthogonal projector onto $\ker D$, and
$Q=I-P$.  Suppose $PA_0P$ is invertible on $P\mathcal H$.  Then
\begin{equation}
p^\dagger(\Gamma D+A_0)^{-1}c
=p^\dagger P(PA_0P)^{-1}Pc+O(\Gamma^{-1}).
\label{eq:singular-D}
\end{equation}
\end{theorem}

\begin{proof}
In the orthogonal decomposition
$\ker D\oplus(\ker D)^\perp$, the Hermitian $D$ has blocks
$0\oplus D_{QQ}$ with $D_{QQ}>0$.  The inverse of
$\Gamma D_{QQ}+A_{QQ}$ is $\Gamma^{-1}D_{QQ}^{-1}+O(\Gamma^{-2})$.
The Schur complement on the $P$ block is therefore
$PA_0P+O(\Gamma^{-1})$, whose inverse is
$(PA_0P)^{-1}+O(\Gamma^{-1})$.  Solving the block system yields
\eqref{eq:singular-D}.
\end{proof}

\begin{scopebox}[Restriction of Theorem~\ref{thm:singular-D}]
The orthogonal decomposition used above requires $D=D^\dagger\ge0$.  A general
Liouvillian fast block is non-Hermitian and can be nonnormal; then
$\ker D$ need not be orthogonal to $\Ran D$.  The same notation must not be
reused with an orthogonal projector.  A nonnormal extension requires a Riesz
spectral projector, an invariant complement, and a Drazin inverse, together
with a semisimplicity/spectral-gap hypothesis.  That extension is not claimed
by Theorem~\ref{thm:singular-D}.
\end{scopebox}

% ==================================================================
\section{Regular domains and singular spectral points}
% ==================================================================

\begin{theorem}[Adjugate identity on analytic parameter domains]
Let $A(\theta)$ be analytic on a connected domain $U$ and define
\[
N(\theta)=p^\dagger\adj A(\theta)c,
\qquad
\Delta(\theta)=\det A(\theta).
\]
On the regular set, $K=N/\Delta$.  If $N\equiv0$, then $K=0$ at every regular
point.  Conversely, if $K$ vanishes on a nonempty open subset of the regular
set, then $N\equiv0$ on $U$.
\end{theorem}

\begin{proof}
The minors defining $N$ and $\Delta$ are analytic.  On the regular set the
adjugate formula gives $K=N/\Delta$.  If $K=0$ on an open regular subset, then
$N=0$ there; the identity theorem gives $N\equiv0$ on the connected domain.
\end{proof}

\begin{theorem}[Semisimple zero mode and pole cancellation]
Suppose $0$ is a semisimple eigenvalue of $A$, let $P_0$ be its spectral
projector, and let $A^D$ be the group/Drazin inverse.  Then
\[
(A+sI)^{-1}=\frac{P_0}{s}+A^D+O(s),
\]
and
\[
p^\dagger(A+sI)^{-1}c
=\frac{p^\dagger P_0c}{s}+p^\dagger A^Dc+O(s).
\]
The pole is removable iff $p^\dagger P_0c=0$.
\end{theorem}

\begin{proof}
Use the invariant direct sum of the zero eigenspace and the complementary
range.  On the first block the inverse is $s^{-1}I$; on the second it is the
regular expansion of an invertible matrix.
\end{proof}

\begin{proposition}[Non-semisimple zero mode]
Let $P_0$ be the Riesz projector onto the generalized zero eigenspace, and let
$N:=AP_0$ be nilpotent of index $k$, so $N^k=0$.  On that subspace,
\begin{equation}
(sI+N)^{-1}=\sum_{j=0}^{k-1}(-1)^jN^j s^{-j-1}.
\label{eq:jordan-principal}
\end{equation}
Therefore the transfer principal part contains the coefficients
\[
(-1)^j p^\dagger N^jP_0c,
\qquad j=0,\ldots,k-1.
\]
The singularity is removable only if every principal-part coefficient
vanishes.
\end{proposition}

\begin{proof}
Multiply the finite geometric sum in \eqref{eq:jordan-principal} by
$sI+N$; all intermediate terms cancel and the remainder is
$I+(-1)^{k-1}N^ks^{-k}=I$.
\end{proof}

\begin{remark}
A numerical $0/0$ at a singular point must never be reported as a transfer
zero before the Laurent principal part has been checked.
\end{remark}

% ==================================================================
\section{Rate conversion used in the NV model}
% ==================================================================

\begin{proposition}[Symmetric orbital hopping]
For
\[
L_1=\sqrt{k}\ket Y\bra X,
\qquad
L_2=\sqrt{k}\ket X\bra Y,
\]
the population imbalance decays at rate
$\Gamma_{XY}=2k$, while each optical coherence to a ground state decays at
rate $k/2=\Gamma_{XY}/4$.
\end{proposition}

\begin{proof}
The populations satisfy
$\dot p_X=-kp_X+kp_Y$ and $\dot p_Y=kp_X-kp_Y$, hence
$\dot{(p_X-p_Y)}=-2k(p_X-p_Y)$.  For an optical coherence $\bra X\rho\ket g$,
the recycling term vanishes because the jumps are confined to the excited
manifold, while the anticommutator from $L_1^\dagger L_1=k\ket X\bra X$
contributes $-k/2$.  The other coherence is treated symmetrically.
\end{proof}

% ==================================================================
\section{PRL-safe theorem statement}
% ==================================================================

\begin{claimbox}[Central statement supported by the theorems]
For a finite-dimensional fixed-scaling response
$K_\Gamma=p^\dagger(\Gamma D+A_0)^{-1}c$ with invertible block-diagonal $D$,
selection rules can force low path moments to vanish.  The first nonzero
moment determines the asymptotic integer dissipation order.  A sector graph
provides a rigorous lower bound through its source-to-readout distance; the
bound becomes an equality only when the summed shortest-path amplitude is
nonzero.  Observable exponents follow by applying the product/quotient rule to
the complete measured response, not to a single kernel in isolation.
\end{claimbox}

The following stronger sentences are \emph{not} consequences of the present
theorems without additional hypotheses:
\begin{itemize}
\item that graph distance always equals the order, regardless of path
      interference;
\item that a reduced-kernel order is automatically the order of normalized
      contrast;
\item that an asymptotic class is experimentally resolvable in a finite
      engineerable window;
\item that the Hermitian singular-damping theorem applies unchanged to a
      nonnormal Liouvillian;
\item that a simplified-model exact zero is an exact zero of the full NV
      Liouvillian.
\end{itemize}

% ==================================================================
\section*{Dependency chart}
% ==================================================================

\begin{center}
\renewcommand{\arraystretch}{1.25}
\begin{tabular}{lll}
\toprule
Result & Main tool & Logical role\\
\midrule
Pure stationary state & purity + exact commutator algebra & closes v6.2 gap\\
Neumann lemma & matrix geometric series & explicit error bound\\
Laurent theorem & rational resolvent at infinity & asymptotic order\\
Sector-graph theorem & noncommutative word expansion & moment lower bound\\
Krylov theorem & minimal polynomial + rational identity & exact zero\\
Inheritance theorem & asymptotic product/quotient & observable order\\
Singular damping & Schur complement & Hermitian $D$ only\\
Jordan proposition & finite nilpotent series & higher-pole test\\
\bottomrule
\end{tabular}
\end{center}

\bigskip
\noindent
The only analysis-level input used above is the identity theorem for analytic
functions.  The large-$\Gamma$ order theorem itself is finite-dimensional
rational-function algebra; the Neumann series is auxiliary rather than
foundational.

\end{document}
